\documentclass{beamer}
\usepackage{hyperref}
\usepackage[T1]{fontenc}

% other packages
\usepackage{latexsym,amsmath,xcolor,multicol,booktabs,calligra}
\usepackage{graphicx,pstricks,listings,stackengine}
\usepackage{multicol}

\title[Reflection on ChatGPT]
{Mid Term Report}
\subtitle{ChatGPT and Academic Integrity}
\author[Eric] % (optional, for multiple authors)
{Li Guanlin\inst{1} \and He Lingrui\inst{1} \and Yan Yuheng \inst{1} \and Li Jiazhou \inst{1}}
\institute[NJU] % (optional)
{
\inst{1}%
Undergraduates in ICS\\
Nanjing University
\and
% \inst{2}%
% Faculty of Chemistry\\
% Very Famous University
}
\date[NJU 2023] % (optional)
{\today}
\usepackage{NJU}

% defs
\def\cmd#1{\texttt{\color{red}\footnotesize $\backslash$#1}}
\def\env#1{\texttt{\color{blue}\footnotesize #1}}
\definecolor{deepblue}{rgb}{0,0,0.5}
\definecolor{deepred}{rgb}{0.6,0,0}
\definecolor{deepgreen}{rgb}{0,0.5,0}
\definecolor{halfgray}{gray}{0.55}

\lstset{
    basicstyle=\ttfamily\small,
    keywordstyle=\bfseries\color{deepblue},
    emphstyle=\ttfamily\color{deepred},    % Custom highlighting style
    stringstyle=\color{deepgreen},
    numbers=left,
    numberstyle=\small\color{halfgray},
    rulesepcolor=\color{red!20!green!20!blue!20},
    frame=shadowbox,
}

\renewcommand{\familydefault}{\rmdefault} %可选择\rmdefault \sfdefault \ttdefault 分别是罗马字体，sans-serif字体和等宽字体  \familydefault是默认字体

\begin{document}
\begin{frame}
    \titlepage
    \begin{figure}[htpb]
        \begin{center}
            \includegraphics[width=0.14\linewidth]{pic/NJU_Logo.eps}
        \end{center}
    \end{figure}
\end{frame}

\begin{frame}
    \tableofcontents[sectionstyle=show,subsectionstyle=show/shaded/hide,subsubsectionstyle=show/shaded/hide]
\end{frame}

\section{Introduction}

\begin{frame}
    \frametitle{Introduction}
    \Large
    \begin{itemize}
        \item Purpose: \textbf{ChatGPT and Academic Integrity.}
        \item Questions: Reasoning and Non-trivial Questions.
              %The capability of ChatGPT in demonstrating its reasoning and answering non-trivial questions.
        \item Siginificance: \alert{The Unprecedented Threat.}%Raising an alerm about the unprecedented threat.
    \end{itemize}
\end{frame}

\section{Methods}

\begin{frame}
    \frametitle{Methods}
    \Large
    \begin{block}{Design}
        Steps:

        \begin{enumerate}
            \item<1-> ChatGPT generates a critical thinking question
            \item<2-> ChatGPT answers the questions itself
            \item<3-> ChatGPT critically evaluates the answer
        \end{enumerate}
    \end{block}
\end{frame}

\begin{frame}
    \frametitle{Evaluation}
    \onslide<1->
    \Large{Standard: Universal intellectual standards}
    \onslide<2->
    \begin{block}{Dimensions}
        \begin{multicols}{2}
            \begin{itemize}
                \item Relevance
                \item Clarity
                \item Accuracy
                \item Precision
                \item Depth
                \item Breadth
                \item Logic
                \item Persuasiveness
                \item Originality
            \end{itemize}
        \end{multicols}
    \end{block}
\end{frame}

\section{Results}
\begin{frame}
    \frametitle{Results}
    \LARGE
    \begin{itemize}[<+->]
        \item Capabilities of ChatGPT : exceptional level
        \item Most significant technological leap
    \end{itemize}
\end{frame}

\begin{frame}
    \frametitle{Descriptive:4 Q\&A in different scenarios}
    \LARGE
    \begin{itemize}
        \item Strengths
        \item Weaknesses
        \item Improvements
    \end{itemize}
\end{frame}

\begin{frame}
    \frametitle{Analytic}
    \begin{block}{Dimesions}
        \begin{multicols}{3}
            \begin{itemize}
                \item<1-> Clarity
                \item<1-> Accuracy
                \item<1-> Precision
                \item<2-> Relevance
                \item<2-> Depth
                \item<2-> Breadth
                \item<3-> Logic
                \item<3-> Persuasiveness
                \item<3-> Originality
            \end{itemize}
        \end{multicols}
    \end{block}
\end{frame}

\section{Discussion}

\subsection{Findings}
\begin{frame}
    \frametitle{Findings}
    \LARGE
    \begin{itemize}
        \item AI's critical thinking
        \item ChatGPT's ability
        \item Impact on exam integrity
    \end{itemize}
\end{frame}

\subsection{Links}
\begin{frame}
    \frametitle{Link to study objective}
    \LARGE
    \begin{itemize}
        \item ChatGPT
        \item Exam Integrity
        \item Threat
    \end{itemize}
\end{frame}

\begin{frame}
    \frametitle{Novelty}
    \LARGE
    \begin{itemize}
        \item Knowing the ability of AI
        \item Raising an alarm about exam integrity
    \end{itemize}
\end{frame}

\subsection{Conclusion}
\begin{frame}
    \frametitle{Conclusion}
    \Large
    \begin{itemize}
        \item<1-> One technology
        \item<2-> One capacity
        \item<3-> One worry
        \item<4-> One final perspective
        \item<5-> Four solutions
    \end{itemize}
\end{frame}

\subsection{Recommendation}

\begin{frame}
    \frametitle{Recommendation}
    \LARGE
    \begin{itemize}[<+->]
        \item Full understanding: a short period of time
        \item Better strategies: mitigating strategies only
    \end{itemize}
\end{frame}

\section{Reference}

\begin{frame}[allowframebreaks]
    \bibliography{ref}
    \bibliographystyle{acm}
    % 如果参考文献太多的话，可以像下面这样调整字体：
    % \tiny\bibliographystyle{alpha}
\end{frame}

\begin{frame}
    \begin{center}
        {\Huge\calligra Thanks!}\cite{origin}
    \end{center}
\end{frame}

\end{document}